\sectionguard
\subsection{Anselm of Canterbury}\label{sec:fa-anselm}

Anselm, monk and abbot of Bec and then archbishop of Canterbury (d.~1109),
whom the Church honours as a Doctor, gives the schools their first whole
treatise on the fall of the angels. \emph{De casu diaboli} is a dialogue
between a master and his disciple, and it begins from a single question of
the Apostle: ``what hast thou that thou hast not received?'' (1~Cor 4:7).
Is it said to angels as well as to men? The master answers that no
creature has anything of itself, and that from the supreme Good there comes
nothing but good and being (c.~1). Then the disciple presses the hardest
form of the difficulty. The angel who stood in the truth persevered
because he had perseverance, had it because he received it, and received it
because God gave it; therefore, it seems, the angel who did not stand did not
persevere because God did not give. Every writer the disciple has read or
heard settles the question by that rule, and he has never seen it answered
(c.~2). The whole dialogue is Anselm's answer.\footnote{\emph{De casu
diaboli} is cited from Migne's reprint (PL~158), by chapter.}

\paragraph{What the angel received.} The master first breaks the rule. When
a man holds out a gift and another takes it, the giving causes the taking;
when he holds it out and it is not taken, the refusal is not caused by his
not giving, but his not giving follows from the refusal (c.~3). The devil
received from God the will and the power to receive perseverance and to keep
the will he had; \latin{sponte dimisit voluntatem quam habebat}, the disciple is
told, and he did not receive the keeping of it because he let it go (c.~3).
The master then shows how one lets go of a thing one holds. A man who holds
a coin and would rather have bread does not give up the coin because he has
ceased to want it; he ceases to want to keep it because he wants something
else that he cannot have without giving it up. So the devil did not desert
justice because he first ceased to will it; he ceased to will it because he
willed another thing, and in willing that thing he let justice go (cc.~3--4).
The master closes the first argument with the sentence that answers the
disciple's rule:

\begin{quote}
\latin{Quamvis igitur bonus angelus ideo accepit perseverantiam, quia Deus
dedit, non ideo tamen malus non accepit, quia Deus non dedit: sed Deus ideo
non dedit, quia ille non accepit; et ideo non accepit, quia accipere
noluit.} (c.~4; PL~158:333)
\end{quote}

\noindent Though the good angel received perseverance because God gave it,
the evil angel did not fail to receive it because God did not give it; God
did not give because he did not receive, and he did not receive because he
would not.

\paragraph{Justice and advantage.} Every rational will, Anselm teaches, can
will only two things: justice, \latin{iustitia}, and advantage,
\latin{commodum}, of which happiness is made (c.~4). To show why God gave
the angel both wills, the master builds an angel in thought, part by part
(cc.~12--14). Suppose the angel made ready to will but not yet willing: he
cannot move himself to his first will, since whatever moves itself to will
must already will to move itself (c.~12). Give him only the will for
happiness: he can will nothing else, and he wills happiness as high as he
can conceive it, even to be like God; yet his will is neither just nor
unjust, because he wills by necessity, and if nothing higher were open to
him he would will the pleasures of the beasts and be no more to blame
(c.~13). Give him only the will for what is right: the same follows, for a
rectitude he cannot help is not justice (c.~14). God therefore gave both
together, so that the angel should will to be happy and should will it
justly,

\begin{quote}
\latin{quatenus addita justitia sic temperet voluntatem beatitudinis, et
resecet voluntatis excessum, et excedendi non amputet potestatem.}
(c.~14; PL~158:346)
\end{quote}

\noindent Justice is added in order that it may temper the will for
happiness: that it may cut back the will's excess and yet not cut off the
power to exceed. Because he wills to be happy, the angel can go beyond the
measure; because he wills justly, he does not will to. The same chapter
draws the end in view: the angel who does not will what he ought not, though
he could, deserves never to be able to will it, and the angel who deserts
justice by an immoderate will deserves to lack every good (c.~14). Justice,
the disciple then agrees, is \latin{aliquid}, a real thing, and a very good one; injustice is
nothing but the absence of the justice that ought to be there, as the want
of a beard is no fault in a boy and a disfigurement in a man (cc.~15--16).
The will that has deserted justice keeps the debt of it, and the debt is
the trace of a nobility received: to have owed and still owe justice shows
the dignity of the nature; not to have it convicts the person (c.~16).

\paragraph{The sin.} What, then, did the devil will? He could not sin by
willing justice, and whatever he already had he was bound to will. He
therefore sinned by willing an advantage he did not yet have:

\begin{quote}
\latin{Peccavit ergo volendo aliquod commodum, quod nec habebat, nec tunc
velle debuit, quod tamen ad augmentum beatitudinis esse illi poterat.}
(c.~4; PL~158:333)
\end{quote}

\noindent He sinned by willing some advantage which he did not have and
ought not then to have willed, though it could have been an increase of his
happiness. The good he reached for was a real good; the sin was the
disorder of reaching for it against the will of God and before God willed
to give it, as Eve willed to be like the gods before God willed it (c.~4).
In that act the devil willed to be like God inordinately:

\begin{quote}
\latin{Etiamsi noluit omnino par esse Deo, sed aliquid minus Deo, contra
voluntatem Dei, hoc ipso voluit inordinate similis esse Deo, quia propria
voluntate, quae nulli subdita fuit, voluit aliquid. Solius enim Dei esse
debet, sic voluntate propria velle aliquid, ut superiorem non sequatur
voluntatem.} (c.~4; PL~158:333)
\end{quote}

\noindent Even if he did not will to be altogether equal to God, but
something less than God against God's will, by that very act he willed
inordinately to be like God, because he willed something by his own will,
subject to none; for to will by one's own will so as to follow no higher
will belongs to God alone. Anselm draws the last consequence at once: in
presuming to have a will of his own the devil made himself equal to God,
and in willing what God did not will him to will he made himself greater,
\latin{quoniam voluntatem suam supra voluntatem Dei posuit}: he set his own
will above the will of God (c.~4). What the advantage was, Anselm does not
name: \latin{Quid illud fuerit non video}, the master says; it is enough to
know that it was something to which the angels could grow, which they did
not receive when they were made, so that they might advance to it by merit
(c.~6). Aquinas reads the chapter as teaching that the devil ``sought that
to which he would have come had he stood fast,'' and finds his own account
of the devil's likeness to God in agreement with it (\ST{I q.~63 a.~3}; see
\S\ref{sec:q63}).

\paragraph{The confirmation of the good and the fall of the evil.} Before
the fall the good angels could sin as the others did; had they been unable,
they would have kept justice by necessity, would have merited nothing, and
could not rightly be called just (c.~5). The two companies divided in one
choice. Those who willed the justice they had rather than the ``more'' they
had not lost, as far as their will went, the good they gave up for
justice's sake, and received it from justice as a reward:

\begin{quote}
\latin{Quapropter adeo sunt provecti, ut sint adepti quidquid potuerunt
velle; nec jam videant quid plus velle possint: et propter hoc peccare
nequeunt.} (c.~6; PL~158:334)
\end{quote}

\noindent They have been so advanced that they have attained whatever they
could will, and see nothing more that they could will; and for this reason
they cannot sin. Those who preferred the ``more'' that God did not yet will
to give them lost it and lost the good they held, \latin{ut adhaerentes
justitiae nullum velle bonum possint, quo non gaudeant; et deserentes illam,
nullum velle queant, quo non careant}: those who cleave to justice can will
no good they do not rejoice in, and those who deserted it can will none they
do not lack (c.~6). Anselm sets the two final states against each other as
punishment and reward:

\begin{quote}
\latin{Palam igitur est quia sicut illi non posse recuperare, quod deseruit,
est poena peccati, ita huic non posse deserere quod tenuit, est praemium
justitiae.} (c.~24; PL~158:357)
\end{quote}

\noindent As the one's inability to recover what he deserted is the penalty
of sin, so the other's inability to desert what he held is the reward of
justice. The good angel's knowledge that sin is punished, which he now has
from the example of the fallen, is not the cause of his steadfastness: had
neither sinned, God would have given both the same knowledge in another way,
as the reward of perseverance; he taught it by the devil's fall only because
his greater power brings good out of evil, so that no disorder remains in
the kingdom of his wisdom (c.~25).

The fallen angel cannot return. Asked whether the deserter can come back to
the will for justice by himself, the disciple answers: \latin{Sed modo multo
minus: tunc enim conditione naturae non poterat habere; nunc vero merito
quoque culpae non debet habere} (c.~17; PL~158:349). Now much less; before,
by the condition of his nature, he could not have it of himself; now, by the
desert of his fault, he ought not to have it. And the gratitude owed to God
is the same on both sides, since both received from him to have justice, to
be able to keep it, and to be able to desert it; the last was given so that
they might in some way give justice to themselves by not taking it away.
\latin{Semper igitur debet Deo gratias agere malus angelus pro beatitudine
quam sibi abstulit; sicut bonus pro ea quam sibi ipse dedit} (c.~18;
PL~158:350): the evil angel owes God thanks for the happiness he took from
himself as much as the good angel for the happiness he gave himself. God
makes the evil angel unjust only by not restoring the justice he could
restore (c.~18).

\paragraph{Neither knew what would follow.} The angel who fell could not
foreknow his fall. Had he foreseen it and willed it, he had already fallen;
had he foreseen it and not willed it, he would have been miserable in
proportion to his justice, which is unfitting before any sin (c.~21). He
knew that he ought not to will what he willed and that he would deserve
punishment if he did, but he did not know that the punishment would follow.
Anselm adds a consideration that looks forward to \emph{Cur Deus homo}: the
angel could not easily have believed that God would damn a creature made in
such goodness, since no example of avenging justice had yet been seen, and he
knew the number of those made to enjoy God to be fixed with such wisdom that
it had nothing superfluous and would be imperfect if diminished (c.~23). The
good angel did not know it either; had he known it, he would have had two
reasons for not sinning, the love of justice, honourable and sufficient, and
the fear of punishment, neither honourable nor needed, and his perseverance
shines the more because love of justice alone was its reason (c.~24).

\paragraph{The will as its own cause.} The dialogue ends by asking why the
angel willed what he ought not. The will to do so, as a power and as an act,
was something, and so was good and from God, who permitted it; the evil in
it was the injustice, which is nothing and is the sinner's own (cc.~19--20,
28). Beyond that there is no cause to find:

\begin{quote}
\latin{Nonnisi quia voluit. Nam haec voluntas nullam aliam habuit causam,
qua impelleretur aliquatenus, aut attraheretur; sed ipsa sibi efficiens
causa fuit, si dici potest, et effectus.} (c.~27; PL~158:360)
\end{quote}

\noindent Only because he willed. This will had no other cause by which it
was driven or drawn; it was, if it can be so said, its own efficient cause
and its own effect. The Fourth Lateran Council confesses the same in the
Church's words: \latin{Diabolus enim et alii daemones a Deo quidem natura
creati sunt boni, sed ipsi per se facti sunt mali} (DS~800).

\paragraph{The places of the fallen (\emph{Cur Deus homo} I.16--18).} In the
first book of \emph{Cur Deus homo} Boso asks Anselm for the reason of
something ``a part of our belief'': that God meant to make up the number of
the fallen angels from human nature. Anselm answers:

\begin{quote}
There is no question that intelligent nature, which finds its happiness,
both now and forever, in the contemplation of God, was foreseen by him in a
certain reasonable and complete number, so that there would be an unfitness
in its being either less or greater. (I.16)
\end{quote}

\noindent Either the fallen angels were made within that number, or they
fell of necessity because they were outside it, which is absurd; the number
must therefore be restored, and it can be restored only from men (I.16).
Other angels cannot be substituted, for they ``ought to be such as the
former angels would have been, had they never sinned'': they would have
persevered without seeing the punishment of sin, which no angel created
after the fall could do. ``For it must not for a moment be supposed that
good angels are upheld by the fall of evil angels, but by their own virtue''
(I.17). Anselm then argues at length that the elect will be more than the
fallen angels, so that no saint need rejoice at another's ruin as the
condition of his own glory, and that man was made for his own sake and not
only as a replacement: ``even had no angel fallen, men would yet have had
their place in the celestial kingdom'' (I.18). Reading Moses' canticle in
both its translations, ``according to the number of the children of Israel''
and ``according to the number of the angels of God'' (Deut 32:8), he
concludes ``that as many men will be taken as there were angels who remained
steadfast'' (I.18). He offers the whole as his opinion, received ``no
further certainty than as my opinion for the present, until God makes some
clearer revelation to me'' (I.18).\footnote{The English of \emph{Cur Deus
homo} is Sidney Norton Deane's (Chicago, 1903).} Aquinas holds, with the text
Hugh of Saint Victor also cites, that more angels stood than fell
(\ST{I q.~63 a.~9}), and that men are taken up into the angelic orders by
grace (\ST{I q.~108 a.~8}).

\sectionguard
\subsection{Bernard of Clairvaux}\label{sec:fa-bernard}

Bernard, abbot of Clairvaux (d.~1153), the Mellifluous Doctor, speaks of the
angels as a monk speaks of the company he expects to keep. He treats them
in four places: in the fifth book of \emph{De consideratione}, written for
his former monk Pope Eugenius~III, where the angels are the first object of
the ascent above the Pope's cares; in his Lenten sermons on the ninetieth
psalm, \emph{Qui habitat}, where the eleventh and twelfth expound ``He hath
given his angels charge over thee'' (Ps 90[91]:11); in his two sermons for
the feast of Saint Michael; and throughout the sermons on the Song of
Songs.\footnote{The English of \emph{De consideratione} is George Lewis's (Oxford,
1908), of the sermons for the seasons that of a priest of Mount Melleray
(Dublin, 1921--1925), and of the sermons on the Song Samuel J. Eales's
(London, 1896); the Latin is from Migne (PL~182--183).}

\paragraph{Faith, understanding, and opinion.} Before turning to the angels
Bernard distinguishes three ways of knowing invisible things: ``Faith is, by
the exercise of the will, a sure foretaste of truth not yet manifested.
Understanding is the sure and clear knowledge of some invisible thing.
Opinion is the holding something provisionally true which you do not know
to be false'' (\emph{De consid.}\ V.3.6). He then directs consideration to
``our mother, Jerusalem above'' and sorts what is known of her citizens by
these three ways:

\begin{quote}
And first let us remember that the citizens of that country are spirits,
mighty, glorious, blessed, separate personalities, of graduated rank, from
the beginning standing in their own order, perfect of their kind, having
ethereal bodies, endowed with immortality, passionless; not so created, but
so made---that is, through grace, not by nature; beings of pure mind,
benignant affections, religious and devout; of unblemished morality;
inseparably one in heart and mind, blessed with unbroken peace, God's
building, dedicated to the divine praises and service. All this we
ascertain by reading, and hold by faith. (V.4.7)
\end{quote}

\noindent Of their bodies, ``some authorities hesitate to say not only
whence they are derived, but whether in any real sense they exist at all,''
and Bernard is content to leave the matter among opinions. That the angels
have understanding is neither of faith nor of opinion, but a conclusion of
the understanding itself, ``for if they had not understanding, they could
not be partakers of the Divine nature'' (V.4.7). Their names are known to us
by hearing, and by the names ``the duties, merits, ranks, orders, of these
blessed ones \dots\ in one way or another may be conjectured and
distinguished'' (V.4.7). Bernard's angelology declares each thing at the
degree at which it is known: the angels' nature and blessedness by faith,
their intelligence by reason, the offices signified by their names by a
reverent conjecture.

\paragraph{The nine names (V.4.8).} The conjecture follows at once, each
clause opening with \latin{Putemus}, ``let us suppose''. Bernard goes up the
nine names in Gregory's sequence. The Angels are those ``who are believed to
have been given as guardians of individual men,'' of whom the Saviour said,
``Their angels do always behold the face of your Father''; the
Archangels, over them, are sent only ``for particular and very weighty
reasons,'' as ``the great Archangel Gabriel was, as we read, sent to Mary,
and for the greatest of all reasons.'' The Virtues are those through whom
signs appear in the elements for the warning of men, and Bernard finds them
in the Gospel, where after ``there shall be signs in the sun, and the moon,
and the stars'' comes ``for the Virtues of the heavens shall be moved'';
the Powers restrain the power of darkness and the malignant spirits of the
air; the Principalities establish, rule, and change all sovereignty on
earth; the Dominions tower above the rest, and to them the lower orders
render account ``as it were to their masters'' (V.4.8; Luke 21:25--26). Of
the Thrones he writes: \latin{Putemus Thronos alto etiam ab his evolasse
recessu, qui ex eo quod sedent, Throni dicuntur; et ex eo sedent, quod sedit
in eis Deus} (V.4.8; PL~182:792): let us suppose that the Thrones have
flown to a height far withdrawn even from these, that they are called
Thrones because they sit, and that they sit because God is seated in them.
Their sitting is ``the deepest tranquillity, the utmost calmness and
serenity, the peace which passes all understanding.'' The Cherubim drink at
the fount of wisdom, ``the mouth of the Most High,'' and pour its streams
upon the city; the Seraphim, set on fire by the divine fire, kindle all
things, ``so that the citizens may be each a burning and a shining lamp,
burning with charity, shining with knowledge'' (V.4.8). The same offices are
expounded under the treatise's order of the orders at \S\ref{sec:q108}.

\paragraph{The orders as mirrors of God (V.4.9--10).} Bernard does not stay
with the orders. ``How good it is, Eugenius, for us to be here! But how much
better will it be if we ever altogether follow on whither we have in part
gone before!'' (V.4.9). The soul that gathered all its affections would
travel round ``the abodes of light'' and see the heart of God revealed there
(V.4.9). Then each order becomes a window on its Maker. In the Seraphim is
seen how he loves who has no cause to love; in the Cherubim that the Lord is
a God of knowledge; in the Thrones the Judge who will not deceive and cannot
be deceived; and so down the scale: ``In the Principalities we may perceive
the fount of all things; as a door turns on its hinge, so the universe
depends on the King Himself. In the Powers we may see how powerfully the
First Cause protects those over whom He rules, keeping off and driving back
the hostile powers'' (V.4.10). The last are nearest to us:

\begin{quote}
Lastly, as we contemplate the Angels and Archangels, we may see, and
marvel, how true it is in our experience that `He careth for us'; He who
never ceases to delight us with the visits of such glorious beings, to
instruct us with their revelations, admonish us through their suggestions,
solace us by their zealous attention. (V.4.10)
\end{quote}

\paragraph{The angel with the soul, God in the soul (V.5.11--12).} All the
gifts of the orders come from ``one and the self-same Spirit dividing to
them severally as He willed'': the Seraphim burn, but with the fire of God;
the Cherubim shine, but by participation in the truth (V.5.11). \latin{Adsunt
Angeli et Archangeli: sed ille germanior nobis, qui non modo adest, sed
inest} (V.5.11; PL~182:795): Angels and Archangels are present with us, but
he is closer kin to us who is not only present but within. Eugenius may
object that an angel can also be within, since Zacharias heard ``the angel
that spoke in me'' (Zach 1:14). Bernard grants it and draws the line:

\begin{quote}
\latin{Inest Angelus suggerens bona, non ingerens: inest hortans ad bonum,
non bonum creans.} \dots\ \latin{Angelus ergo cum anima, Deus in anima.
Ille ut contubernalis animae inest, Deus ut vita.} (V.5.12; PL~182:795)
\end{quote}

\noindent The angel is within as one who suggests good things, not as one
who puts them in; as one who urges to the good, not as one who creates it.
``The Angel, therefore, is with the soul, God is in the soul. The Angel is
in the soul as a comrade, God as life'' (V.5.12, Lewis). God loves like
Charity, knows like Truth, judges like Equity, rules like Majesty, guards
like Safety, works like Virtue, reveals like Light; ``All these things the
Angels also do, and so do we, but in a far inferior way, not, of course, by
our native goodness, but by the goodness whereof we partake'' (V.5.12).
Aquinas determines the same boundary: an angel can persuade the will by
presenting a good to the intellect, but God alone moves it from within
(\ST{I q.~106 a.~2}; \ST{I q.~111 a.~2}).

\paragraph{How each order loves (\emph{In Cant.}\ 19).} Expounding
``therefore young maidens have loved thee'' (Cant 1:2[3]), Bernard asks why each
of the heavenly orders loves the Lord Jesus, and finds the answer in what
each contemplates. The Angels rejoice to execute the judgments of God; the
Archangels are admitted more nearly into the counsels of the eternal Wisdom;
the Virtues search out ``the secret and permanent causes of miracles'';
the Powers magnify ``the Almighty power of our Crucified Lord'';
the Principalities see that he is ``the Firstbegotten of all Creation, and
the Principle of existence to the whole universe''; the Dominations trace
the dominion of Christ that no resistance can check (19.2--3). Upon the
Thrones God is seated, and from them, ``as from a solemn auditorium,'' he
teaches knowledge to angels and to men (19.4). The Cherubim drink from the
Fountain himself, the Lord Jesus, and the Seraphim are so drawn into God,
who is Love, that they seem one spirit with him, ``just as the fire which
enkindles air, and impresses upon it heat and colour like its own, so that
it is manifestly not merely ignited, but itself become fire'' (19.5). Then
Bernard gathers the whole choir into one sentence:

\begin{quote}
God is, then, loved by the Angels because of the supreme righteousness of
His judgments; by the Archangels because of the supreme wisdom of His
purposes; by the Virtues on account of the most blessed showing forth of
miracles and wonders by which He deigns to draw the unbelieving to faith;
but by the Powers on account of that power, equally just and prevailing, by
which He is wont to defeat the cruelty of the wicked, and to defend the good
against them; by the Principalities on account of that eternal and
primordial power which bestows the principle of existence, and existence
itself, upon every creature \dots\ and, lastly, by the Séraphin, inasmuch as
He is Himself Love, that He hateth nothing that He has made, and wills that
all men should come to the knowledge of the Truth and be saved. (19.6)
\end{quote}

\noindent ``All these spirits, then, love God according to the degree of
knowledge they have of Him'' (19.7). The sermon keeps Gregory's sequence of
the orders, as \emph{De consideratione} does.

\paragraph{The angels and bodies (\emph{In Cant.}\ 5).} The fifth sermon
distinguishes four orders of spirits, the beast's, man's, the angel's, and
God's, and asks what need each has of a body. The heavenly spirits need
bodies not for themselves but for their ministry: ``how can they without
bodies fulfil their ministry, especially towards those who are living in
bodies?'' They were seen by the patriarchs, entered their tents, ate with
them, and washed their feet (5.2). For their own knowledge they need none,
since the heavenly spirit reaches the highest things ``by its native force
and the kinship of its nature to them'' and reads without contradiction in
the Book of Life (5.4). Whether their bodies are natural to them, or are
taken and laid aside for their visible appearances, Bernard declines to
determine, since ``the Fathers appear to have held diverse opinions''
(5.7). One thing he teaches without reserve: ``no created spirit can by
itself reach unto our minds \dots\ This prerogative is reserved for the
supreme Spirit'' (5.8). Aquinas cites a sentence of the sixth sermon for the
opinion that angels have bodies naturally united to them and reads it of
assumed bodies (\ST{I q.~51 a.~1 obj.~1 and ad~1}); the Summa determines
that the angels are incorporeal and assume bodies for their ministry
(\ST{I q.~51 aa.~1--2}; see \S\ref{sec:q51}).

\paragraph{The angels at the Office and the one Bride (\emph{In Cant.}\ 7,
27).} The friends of the Bridegroom through whom the Bride asks for his
kiss are the holy angels, ``who are present with human beings who pray, and
offer to God their prayers and vows,'' as Raphael told Tobias that he had
offered his prayer to the Lord (\emph{In Cant.}\ 7.4; Tob 12:12). They
associate themselves with those who sing psalms, as the psalm says, ``I will sing praise to thee in the sight of the angels''
(Ps 137[138]:1).
Bernard rebukes the monks who sleep at the night Office: they do not honour
the citizens of heaven, who are touched by their vigilance ``and delight to
assist at your solemn services'' (7.4). In the twenty-seventh sermon the
whole heavenly company is the Bride's pattern. God ``has created some
Angels, others Archangels, but others Virtues, Dominations, Principalities,
Powers, Thrones; others, again, Chérubin, and others Séraphin. These are the
stars with which His heaven is studded''; and they are so benign that
``some of them do not think it burdensome to remain among us, to be actively
employed in watching over us and doing us good'' (27.5). The Church and the
angels are one Bride:

\begin{quote}
But just as He desired out of diverse flocks of sheep to make one, so that
there should be one Fold, and one Shepherd (John x.~16); so, having from the beginning of
the world one Bride, with whom He was closely united---I mean the multitude
of the angels---it pleased Him to call together a Church out of men also,
and with this to unite that which is from heaven, so that there shall be
one Bridegroom and one Bride. (27.6)
\end{quote}

\paragraph{The ways of the angels (\emph{In Ps.\ Qui habitat} 11).}
Commenting ``to keep thee in all thy ways,'' Bernard examines four ways: of
men, of the demons, of the holy angels, and of the Lord. The ways of the
demons are presumption and obstinacy. Lucifer began with presumption, ``I
will exalt my throne above the stars of God \dots\ I will be like the Most
High'' (Isa 14:13--14), and ``Presumption has made each of them `a spirit
that goeth,' and obstinacy `a spirit that returneth no more''' (11.4). The
ways of the holy angels are those the Only-Begotten revealed when he said
that the disciples would see the angels ascending and descending upon the
Son of man (John 1:51):

\begin{quote}
Therefore the ways of the angels are ascents and descents. They ascend for
their own sakes; they descend or rather condescend to us and for us. For
these blessed spirits mount up to God in blissful contemplation, and out of
compassion for us they come down to our lowliness to keep us in all our
ways. They ascend to rejoice in God's presence; they descend to execute His
will, ``for He hath given His angels charge over thee.'' Nevertheless, not
even when they descend to us are they deprived of the vision of glory,
because they ``always see the Face of the Father.'' (11.6)
\end{quote}

\noindent Aquinas gives the reason in the terms of the schools: the angels
order their outward actions by intellect alone, and their ministry does not
impede their contemplation (\ST{I q.~112 a.~1 ad~3}).

\paragraph{Reverence, devotion, confidence (\emph{In Ps.\ Qui habitat} 12).}
The twelfth sermon is Bernard's great teaching on the guardian angels. The
devil himself, by quoting this psalm against the Lord in the desert, gave
the Church the occasion to sing it through Lent (12.3). God has sent his
Son, given his Spirit, promised the vision of his face, ``And lest there
should be any one at all in the kingdom of heaven unoccupied in the care of
us, Thou hast appointed the blessed angels `to minister for us,' Thou hast
charged them with our keeping, and ordered them to act as our guides''
(12.3). Who gave the charge? The Divine Majesty. To whom? ``to those sublime
and blissful spirits, who dwell so near Himself \dots\ and are truly the
`domestics of God.''' Concerning whom? ``Lord, `what is man that Thou art
mindful of him?''' (12.4). Then comes the heart of the sermon:\footnote{The Mount Melleray translator notes
(1921) that these paragraphs are read in the Roman and Cistercian breviaries
as the lessons of the second nocturn of the feast of 2 October; on the feast
see \S\ref{sec:liturgy}.}

\begin{quote}
\latin{Quantam tibi debet hoc verbum inferre reverentiam, afferre
devotionem, conferre fiduciam! reverentiam pro praesentia, devotionem pro
benevolentia, fiduciam pro custodia. Caute ambula, ut videlicet cui adsunt
angeli, sicut eis mandatum est, in omnibus viis tuis. In quovis diversorio,
in quovis angulo, angelo tuo reverentiam habe. Tunc audeas illo praesente,
quod vidente me non auderes?} (12.6; PL~183:233)
\end{quote}

\noindent How much reverence this word should bring you, how much devotion,
how much confidence: reverence for their presence, devotion for their
goodwill, confidence for their guardianship. Walk warily, as one with whom
the angels are present in all his ways, as they have been commanded. In
whatever lodging, in whatever corner, have reverence for your angel. Would
you dare in his presence what you would not dare with me watching? The
Mount Melleray translator renders it: ``In every place, whether public or
private, show respect for thy angel. Surely thou wouldst not dare to do in
his presence what thou wouldst fear to do in mine'' (12.6). If one doubts a
presence he cannot see, faith, which the Apostle calls ``the evidence (argumentum)
of things that appear not,'' proves it: ``The angels are present, therefore,
and they are present to thee: not only are they with thee, but they are
with thee to defend thee'' (12.6; Heb 11:1).

Honour goes first to God, who gave the charge; yet ``we must also show
ourselves grateful to these celestial spirits, who obey the divine
injunction with so much charity, and assist us in our so great necessity.
\dots\ Let us give them love for love'' (12.7). Bernard explains the
Apostle's ``to God alone be honour and glory'' by the sun, whose light does
not extinguish the stars but covers them: the honour of God prevails in all
honour given to others, so that he ``is honoured in all others,'' and ``let
us in Him love His angels, and love them tenderly, since we are destined as
`joint-heirs' with them to share in their glory'' (12.7). Why fear under
such guardians? \latin{Fideles sunt, prudentes sunt, potentes sunt: quid
trepidamus?} (12.8; PL~183:234): they are faithful, prudent, and mighty;
why do we tremble? The ``hands'' in which they bear us up (Ps 90[91]:12)
are spiritual helps; Bernard names two: the thought of the brevity of
tribulation and the thought of the eternity of the reward, printed on the
heart, ``since it is certain that evil suggestions come from the bad
angels'' (12.10). The sermon ends with a counsel: \latin{Habetote familiares
angelos, fratres mei, frequentate eos sedula cogitatione et devota oratione,
qui semper vobis adsunt ad custodiam et consolationem} (12.10;
PL~183:235): make the angels your familiar friends, my brothers; visit them
often by attentive thought and devout prayer, for they are always with you
to guard and console. The next sermon adds that the same holy angels ``will
bear us up in their hands when life's journey is ended,'' as Lazarus ``was
carried by angels into Abraham's bosom'' (13.1; Luke 16:22). Aquinas takes
the same verse of the psalm for the \emph{sed contra} of the question on
guardianship (\ST{I q.~113 a.~1}; \S\ref{sec:guardians}).

\paragraph{Why the angels love us (\emph{In festo S.\ Michaelis} 1--2).} On
Michaelmas Bernard confesses that he cannot speak of the angels' glory, and
speaks instead of their kindness, since ``these exalted creatures show
themselves not less amiable by their condescension than they are admirable
by their state'' (1.1). They are ministering spirits (Heb 1:14), and no
angel disdains ministry, since the King of angels came ``not to be
ministered unto but to minister'' (Matt 20:28). Christ's ministry surpasses
theirs: ``The angels minister indeed, yet not of their own. They offer to
the Lord for us, not their own, but our good works, and bring us back in
return His graces''; hence the angel of the Apocalypse offers the prayers
of the saints only after ``there was given to him much incense'' (1.2; Apoc
8:3). Christ offered himself, and ministers his own flesh to us ``even to
this day'' (1.3). The angels love us ``because we have been loved so much by
Christ'': ``We, O happy spirits, we are the poor little dogs of that mighty
Lord Whom you love with such ardent affection'' (1.3). Bernard gives three
further reasons for their care, and names them in one Latin sentence:

\begin{quote}
\latin{Hic est, dilectissimi, funiculus triplex, quo de excelso caelorum
habitaculo ad consolandos, ad visitandos, ad adjuvandos nos attrahitur
supereminens charitas angelorum, propter Deum, propter nos, propter se
ipsos.} (1.4; PL~183:449)
\end{quote}

\noindent This is the threefold cord by which the surpassing charity of the
angels is drawn down from its high dwelling to console, visit, and help us:
for God's sake, for ours, and for their own. For God's sake, because they
imitate his mercy; for ours, because they pity in us a likeness to
themselves, since human souls, capable of reason and of beatitude, are akin
to them; for their own, because they long to see the ruined walls of their
city rebuilt from the living stones of men (1.4). What pleases them most in
us is unity and peace, ``the angels of peace'' finding in concord ``a new
Jerusalem on earth,'' and what drives them away is discord (1.5--6). The
second sermon applies the Gospel of the feast, the Lord's warning against
scandalising the little ones whose angels see the Father's face (Matt
18:6--10): in one who does not love his brother the angels find none of the
three reasons for loving him (2.2).

\sectionguard
\subsection{Hugh of Saint Victor}\label{sec:fa-hugh}

Hugh, canon of the abbey of Saint-Victor at Paris (d.~1141), wrote in
\emph{De sacramentis christianae fidei} the first great ordered account of
the whole faith in the schools, and the fifth part of its first book is a
treatise on the angels: \emph{De creatione angelorum et natura, et
confirmatione et lapsu} (PL~176:245--264). Its thirty-four short chapters set
the agenda that the Lombard took over, often word for word, and through him
the schools. Hugh also commented the \emph{Celestial Hierarchy} in the
Latin of John Scotus (PL~175:923--1154).

\paragraph{When and how they were made (I.5.1--7).} Hugh proposes to ask when
the angels were created, what they were made, what they became when
divided by aversion and conversion, and how they were disposed when judged
by damnation and glorification; then their number, orders, excellence,
gifts, offices, and God's government of them (I.5.1). Scripture says both
that wisdom was created first of all (Eccli 1:4) and that ``In the
beginning God created heaven, and earth'' (Gen 1:1). Hugh reconciles them:
the rational creature was made first not in time but \latin{causa solum et
respectu et dignitate}, in cause, relation, and dignity, since the bodily
creation is ordered to it as it is ordered to God, and it alone was made to
God's likeness (I.5.3). In time the spiritual and the corporeal creatures
began together, \latin{simul in tempore uno et in principio temporis}
(I.5.4), as ``He that liveth for ever created all things together'' (Eccli
18:1). Some reckon ages in which the angels served God before the world;
Hugh prefers the other view \latin{salva reverentia secretorum in quibus
nihil temere asserendum putamus} (I.5.4; PL~176:249): saving the reverence
due to hidden things, in which nothing is to be asserted rashly. The
angelic nature was formed and unformed at once: formed by wisdom in the
habit of its nature, unformed in the form it was yet to receive
\latin{per amorem et conversionem ad creatorem suum} (I.5.5), through love
and conversion to its Creator. It was not made from a pre-existing matter,
but as spirits \latin{personaliter distincta et expressa}, personally distinct and expressed, with innate reason and
free choice, able without violence to turn either way, \latin{in altero
adjuti sine coactione, in altero derelicti sine oppressione} (I.5.7;
PL~176:250): helped without compulsion to the one, left without oppression
to the other.

\paragraph{Four properties and their degrees (I.5.8--14).} Hugh gives the
angels four properties in their first condition: \latin{substantiam
simplicem et immaterialem; secundum, discretionem personalem; et tertium
formam sapientiae et intelligentiae rationalem; quartum vero sive ad bonum
sive ad malum liberam inclinandae voluntatis et electionis propriae
potestatem} (I.5.8; PL~176:250): a simple and immaterial substance, personal
distinctness, the rational form of wisdom and intelligence, and the free
power of inclining the will and its own choice to good or evil. In these
they differed by degree, as bodies differ in purity, from the first moment
of their making: those created more subtle in nature and keener in wisdom
were set higher in dignity, \latin{ita tamen ut ii qui excellentiores
constituti sunt sine superbia praelati essent, et qui inferiores positi
sunt, sine miseria subjecti} (I.5.10; PL~176:251): so that the higher were
set over others without pride and the lower placed beneath them without
misery. Their power was threefold, over themselves, over one another, and
over what was placed beneath them: \latin{Prima potestas virtutis erat,
secunda dominationis, tertia administrationis} (I.5.13), of strength, of
dominion, of administration. Their knowledge was threefold too: they knew
what they had been made, by whom, and with whom, so that they understood
what to seek and what to refuse, their beginning and end, and what they
owed each other (I.5.14). Before these distinctions Hugh confesses the
limits of the inquirer: \latin{Nos vero caligante intelligentia timide
incedimus ad illa, et palpamus sensu cognitionis infirmo quae virtute non
comprehendimus} (I.5.12; PL~176:251): we approach these things fearfully,
with clouded understanding, and feel with a weak sense of knowledge what we
cannot grasp by power.

\paragraph{Perfect in the first degree; not foreknowing (I.5.15--22).} Were
the angels made perfect? Hugh distinguishes three perfections: of the
nature made, of the nature glorified, and of the uncreated nature. The
angels were made perfect in the first, to come to the second under the
third; had they begun in the second, there would have been no room left for
them to grow or to fail (I.5.16--17). They did not foreknow their future:
\latin{Propterea monstratum est eis quid faciendum esset non quid esset
futurum} (I.5.18; PL~176:253): they were shown what they were to do, not
what was to come, so that they might choose freely, directed by reason and
not driven by any necessity of foresight. All were made good, since the
Creator cannot be the author of evil; good, just, and blessed with the
goodness nature had received at its beginning, not with any they had made
or merited (I.5.19).

\paragraph{Conversion and aversion (I.5.23--26).} The angels, one in their
created goodness, began to be divided: some rose above what they had been
made and became good by justice, others fell below it and became evil by
fault. Both movements were of the will, and the will of liberty (I.5.23).
Hugh states the part of grace in each with exactness:

\begin{quote}
\latin{Qui autem ad bonum convertebantur voluntarie movebantur gratia
cooperante sine coactione. Qui autem a bono avertebantur sponte
praecipitabantur gratia deserente sine oppressione.} (I.5.24;
PL~176:257)
\end{quote}

\noindent Those who turned to the good were moved willingly, grace working
with them without compulsion; those who turned from the good fell headlong
of their own accord, grace leaving them without oppression. Those who fell
were not left by grace first and then fell; they were left in the very act
of turning away, as the others were taken up in the very act of turning
(I.5.24). The evil lay neither in the will, nor in its motion, nor in the
thing it moved toward, all of which were from God, but in moving beyond the
measure given, \latin{extra mensuram} (I.5.26). Of the wicked will God is
the orderer and not the maker: \latin{Malarum enim voluntatum ordinator est
Deus, non Creator} (I.5.27; PL~176:258). He restrains the evil angels in
four ways: by the limit of the power given them by nature, by a sudden and
hidden act of his power that makes a thing impossible to them, by the
obstacle of other powers set against them, and by the inward judgment by
which he bends even wills that resist him to fulfil his will (I.5.28). The
teaching is taken up at \S\ref{sec:assaults}.

\paragraph{Nine orders and the tenth (I.5.30--33).} On the orders Hugh
writes:

\begin{quote}
\latin{De ordinibus angelorum hoc auctoritas promulgavit, novem in principio
conditos a Deo ordines angelorum, id est angelos, archangelos, principatus,
potestates, virtutes, dominationes, thronos, cherubim, et seraphim, ut
addito homine denarius consummaretur perfectionis caelestis.} (I.5.30;
PL~176:260)
\end{quote}

\noindent Authority has declared that nine orders of angels were founded by
God in the beginning, the Angels, Archangels, Principalities, Powers,
Virtues, Dominations, Thrones, Cherubim, and Seraphim, so that with man
added the tenfold number of heavenly perfection might be completed. Hugh's
list runs upward in the Dionysian order. Man was not made only because the
angels fell, as some think; but because man, created afterwards, is led to
the place from which they fell, the society diminished by their fall is
repaired by him. Tradition holds that some fell from every order, since the
Apostle names ``principalities and powers'' of darkness (Eph 6:12);
\latin{Nusquam tamen in Scriptura nequam spiritus seraphim appellatos
invenio} (I.5.30; PL~176:261): yet nowhere in Scripture does Hugh find evil
spirits called Seraphim, since whatever other gifts they keep in their
malice, they could not keep charity. On the number of the fallen there is
no certain authority, \latin{Omnino tamen probabile videtur plures
remansisse quam cecidisse} (I.5.31): it seems altogether probable that more
remained than fell, as Eliseus said, ``there are more with us than with
them'' (4~Kgs 6:16), which is the \emph{sed contra} of Aquinas's article
(\ST{I q.~63 a.~9}). The names of the orders \latin{non propter se sed
propter nos data sunt} (I.5.32; PL~176:261): were given not for their sake
but for ours, since those known to themselves by contemplation become known
to us by their names. Each order is named from the gift it holds most
excellently, though in the heavenly country all have all, not equally: the
Seraphim, who excel all in love and in knowledge, are named from the
greater gift, charity, and leave the other names to the orders after them
(I.5.32). Whether all are sent, Hugh reports three opinions: that some go
out and others always stand within (Dan 7:10); that all are sent, some by
office and some by occasion, taking the name of their ministry for the
time; and that the three highest orders always stand before God, the three
lowest are sent, and the three middle pass the divine command from the
higher to the lower, so that all are rightly called angels, messengers
(I.5.33). Aquinas determines the question in \ST{I q.~112 aa.~2--4}.

\paragraph{Love goes in where knowledge stays outside (\emph{In Hier.\
cael.}\ VI).} Expounding the Seraphim of the seventh chapter of the
\emph{Celestial Hierarchy}, whose name, Dionysius says, teaches their
``warmth, and keenness'' (\emph{CH} 7.1), Hugh
writes one of the most quoted sentences of the Victorine school. Love is
sharp because it cannot rest outside the beloved, but must go through
everything until it comes to him, and into him:

\begin{quote}
\latin{Si tamen hoc intelligi potest, quoniam dilectio supereminet
scientiae, et major est intelligentia. Plus enim diligitur, quam
intelligitur, et intrat dilectio, et appropinquat, ubi scientia foris
est.} (\emph{In Hier.\ cael.}\ VI; PL~175:1038)
\end{quote}

\noindent If indeed it can be understood, for love surpasses knowledge and
is greater than understanding. More is loved than is understood, and love
goes in and draws near where knowledge stands outside. Aquinas gives the
reason in his reply on the Seraphim and Cherubim: to love the higher things,
and God above all, is better than to know them (\ST{I q.~108 a.~6 ad~3}).
The Dionysian chapter itself is expounded at \S\ref{sec:ch7}.

\sectionguard
\subsection{Honorius Augustodunensis}\label{sec:fa-honorius}

Honorius, the monk and teacher called Augustodunensis (fl.~early twelfth
century), wrote in the \emph{Elucidarium} a short catechism in questions and
answers between a disciple and his master. Its first book gives the
angels a sequence of plain answers (PL~172:1113--1116). God, \latin{ut praepotens
rex}, as a mighty king, first built himself a palace, the kingdom of heaven, and
predestined to it a number of chosen knights: \latin{quem nec liceret
excedi, et quem necesse esset compleri} (I.6), a number that might not be
exceeded and had to be filled. He set the number at ten, \latin{novem
quidem ordinibus angelorum, et decimo hominum}: nine orders of angels and a
tenth of men. Why nine? \latin{Propter Trinitatem, in novenario enim
numero ternarius tertio fit repetitus}: for the Trinity, since in nine the
three is thrice repeated. Why one of men? For the unity, so that the Unity
in Trinity might be praised by angels and men, the spiritual creature and
the corporeal (I.6). The angels were made when it was said, \latin{Fiat
lux} (Gen 1:3); their nature is \latin{spiritualis ignis}, spiritual fire, as it
is written, \latin{Qui facit angelos suos flammam ignis} (I.6; Heb 1:7).
They have no need of names: \latin{Tanta scientia est in angelis, ut non
indigeant nominibus} (I.6), such is their knowledge; the names Michael,
Gabriel, and Raphael are rather surnames given by men from what they do.

The first angel, seeing himself surpass all the orders in glory, willed to
be equal to God and to rule the others as a tyrant; he was cast from the
palace into prison, and the brightest became the darkest (I.7). Honorius
teaches that he fell at once: \latin{Non plenam horam in veritate stetit,
quia mox ut creatus est cecidit} (I.7; PL~172:1114), he did not stand a
full hour in the truth, lest he should taste anything of the inward
sweetness who so early usurped so great a majesty. The fallen could not
return, since as they fell with none urging them they would have had to
rise with none helping them; and Christ did not redeem them, because the
angels were not born from one angel as men are from one man, and one
redeemed angel would have left the rest outside, while an angel, being
immortal, could not die for satisfaction (I.8). Other angels could not be
created to replace them, for a replacement would have had to be such as
the fallen would have been had they stood without seeing any punishment
(I.8), the argument of Anselm's \emph{Cur Deus homo} I.17. The good were
confirmed at once:

\begin{quote}
\latin{Post lapsum illorum mox ita confirmati sunt, ut nunquam cadere nec
peccare possint.} --- \latin{D. Quid est, non possint?} --- \latin{M.
Nunquam velint.} (I.10; PL~172:1115)
\end{quote}

\noindent After the others fell they were at once so confirmed that they can
never fall or sin. --- What does ``cannot'' mean? --- That they will never
will to. Their confirmation came not from the others' fall but from their
own merit: \latin{Cum enim viderent illos malum superbiendo eligere,
indignati sunt, et summo bono fortiter inhaeserunt}, when they saw the
others choose evil in their pride, they were indignant and clung strongly
to the highest Good (I.10; PL~172:1116). Their form is in a manner God's:
\latin{Ut enim imago cerae imprimitur signaculo, sic expressa est in eis Dei
similitudo}, as an image is pressed into wax by a seal, so the likeness of
God is pressed into them, in that they are light, incorporeal, and adorned
with every beauty; and \latin{Nihil est in rerum natura quod eos lateat, cum
in Deo omnia conspiciant}: nothing in nature is hidden from them, since
they see all things in God (I.10). Man was made tenth, \latin{ut
compleretur electorum numerus} (I.11), so that the number of the elect
might be filled. The demons know much from their angelic nature but not the
future, nor thoughts, except as God reveals them, and they can do evil
\latin{non tamen quantum volunt, sed quantum a bonis angelis
permittuntur} (I.9; PL~172:1115): not as much as they will, but as much as
the good angels permit.

The second book gives the guardian angels a chapter:

\begin{quote}
\latin{D. Habent homines custodes angelos? --- M. Unicuique genti,
unicuique civitati praesunt Angeli, qui jura, leges, mores juste dispensant
et ordinant. Unaquaeque etiam anima, dum in corpus mittitur, angelo
committitur, qui eam semper ad bonum incitet, et omnia opera ejus Deo et
angelis in caelis referat.} (II.28; PL~172:1154)
\end{quote}

\noindent Angels preside over every nation and every city, justly
dispensing and ordering their laws and customs; and every soul, when it is
sent into the body, is entrusted to an angel, who always urges it to good
and reports all its works to God and the angels in heaven. God and the
angels, who see all things in God, need no report; the angel's reporting
is his rejoicing at our progress, as there is joy before the angels over
one sinner who repents (Luke 15:10), and his grief is his indignation at our
sins. The angels come to help when there is need, and most of all when they
are invited by prayer, for they pass from heaven to earth in a moment, and
wherever they are sent they always see the face of the Father. They appear
in human form, in a body taken from the air, visible rather than palpable,
and seen only by those to whom they will to show themselves (II.28). On the
guardians in the treatise see \S\ref{sec:guardians}; on assumed bodies,
\S\ref{sec:q51}.

\sectionguard
\subsection{Hildegard of Bingen}\label{sec:fa-hildegard}

Hildegard, abbess of Rupertsberg near Bingen (d.~1179), whom Benedict~XVI
proclaimed a Doctor of the Church in 2012, received the visions of the
\emph{Scivias} and wrote them down with their interpretation, each vision
told and then explained by the heavenly voice. The sixth vision of the
first book is the vision of the choirs of angels (PL~197:437--441).

\paragraph{The vision.} The vision opens:

\begin{quote}
\latin{Deinde vidi in altitudine caelestium secretorum duas acies supernorum
spirituum multa claritate fulgentes, et qui in prima acie erant velut pennas
in pectoribus suis habebant et facies ut facies hominum prae se ferebant, in
quibus et vultus hominum quasi in pura aqua apparebant.} (Scivias I.6;
PL~197:437)
\end{quote}

\noindent Then I saw, in the height of the heavenly secrets, two ranks of
spirits shining with great brightness; those in the first rank had, as it
were, wings on their breasts and faces like the faces of men, in which the
faces of men appeared as in clear water. In the second rank the image of the
Son of man shone in the faces as in a mirror. These two ranks were ringed,
like a crown, by five more: the first of the five shone from the shoulders
down; the second was too bright to look upon; the third was like white
marble, with heads like men's, burning torches above them, and an iron cloud
about them from the shoulders down; the fourth had human faces and feet,
helmets, and tunics of marble; and the fifth showed no human form at all,
\latin{velut aurora rubebant}, glowing red like the dawn. Around these five
two further ranks stood like a crown: the first full of eyes and wings, a
mirror in every eye and a human face in every mirror, their wings raised to
the height above; the second burning like fire, \latin{plurimasque pennas
habentes in quibus quasi in speculo omnes ordines ecclesiasticae
institutionis insignitos demonstrabant} (PL~197:437): with many wings, in
which, as in a mirror, they showed all the orders of the Church's
institution marked out. And all the ranks resounded with every kind of
music the wonders that God works in blessed souls, and glorified God
magnificently (PL~197:437).

\paragraph{The interpretation.} The voice explains that God made some
creatures to cleave to earthly things and others to dwell in the heavenly,
and disposed the blessed spirits both for the salvation of men and for the
honour of his name: some to help men in their needs, others to make known
to men the judgments of his secrets (PL~197:437). The two inner ranks are
the Angels and the Archangels, and they signify that body and soul must
serve God. The Angels' wings are the desires of their deep understanding,
\latin{non quod pennas ut aves habeant, sed quod voluntatem Dei in desideriis
suis velociter perficiant, velut homo in cogitationibus suis celeriter
volat} (PL~197:438): not that they have wings like birds, but that they
swiftly accomplish God's will in their desires, as a man flies swiftly in
his thoughts; and the human faces in them show that they watch God's will
among men and show him men's deeds. The Archangels bear the image of the Son
of man because they \latin{incarnatum Verbum Dei purissime magnificant, quia
ipsi, arcana Dei cognoscentes, mysteria incarnationis Filii Dei signis suis
frequenter praeinsinuabant} (PL~197:438): most purely magnify the incarnate
Word of God, since, knowing the hidden things of God, they often foretold
the mysteries of the Son's Incarnation by their signs. The five ranks that
encircle them are the five senses of man, cleansed by the five wounds of the
Son. The first of the five are the Virtues, who rise into the hearts of
believers and build in them a high tower of charity; the second, too bright
to look upon, are the Powers, for no mortal weakness can grasp the serenity
of God's power; the third, like marble with torches overhead, are the
Principalities, who show that princes of men must put on the firm strength
of justice and look to Christ their head; the fourth, helmed, are the
Dominations, who show that the Lord of all raised the reason of man from
the dust when he sent his Son; the fifth, red like the dawn, are the
Thrones, who show that the Godhead bowed down to humanity when the Only
Begotten took flesh without stain in the dawn, that is, in the blessed
Virgin (PL~197:438--440). The two outer ranks are the Cherubim and the
Seraphim. The Cherubim, full of eyes, \latin{scientiam Dei significantes}
(PL~197:440), signify the knowledge of God, in which they see the mysteries
of the heavenly secrets and foresee those who lift their desires to God as
on wings. The Seraphim burn with the love of God and have the greatest
desire to see him; in their wings they show the Church's orders, secular and
spiritual, because the secrets of God appear wonderfully in them, and all
who seek the life above with a pure heart are to love God ardently and so
come to the joys of those they imitate (PL~197:440). The music of the
choirs is the joy of the blessed spirits in the wonders God works in his
saints, as the Psalmist says, ``The voice of rejoicing and of salvation is
in the tabernacles of the just'' (Ps 117[118]:15; PL~197:441).

Hildegard's vision orders the choirs from within outward, the Angels and
Archangels, then the Virtues, Powers, Principalities, Dominations, and
Thrones, then the Cherubim and the Seraphim. It is Gregory's sequence of the
orders. Its teaching is that of the Fathers, the angels as ministers of
salvation, the names as offices, the Seraphim as love and the Cherubim as
knowledge, given in images, and each order is read as a lesson for the
Church on earth.

\sectionguard
\subsection{Peter Lombard}\label{sec:fa-lombard}

Peter Lombard, master in the schools of Paris and bishop of Paris
(d.~1160), gathered the Church's teaching into the four books of the
\emph{Sentences}, which became the textbook every master of theology
lectured on; Bonaventure, Albert, and Aquinas wrote their commentaries on
it. The Fourth Lateran Council confessed that God \latin{simul ab initio
temporis utramque de nihilo condidit creaturam, spiritualem et corporalem,
angelicam videlicet et mundanam} (DS~800), and in its next chapter professed
the faith in the Trinity \latin{cum Petro Lombardo} (DS~804). The
Master's treatise on the angels fills the second book from the second
distinction to the eleventh, and its order of questions is largely Hugh's,
often in Hugh's words.\footnote{The Latin of the \emph{Sentences} is that
of the Quaracchi edition (1916).}

\paragraph{When, where, and what they were made (dd.~2--4).} The rational
creature is angelic and human, the first wholly spiritual, the second partly
spiritual and partly corporeal, and the treatise begins with the angels so
that reason may pass from the Creator to the worthier creature (d.~1 c.~6).
The first questions are Hugh's: when, where, and what the angels were made;
what they became by aversion and conversion; and their excellence, orders,
gifts, offices, and names (d.~2). On the time, the Master holds
\latin{quod simul creata est spiritualis creatura, id est angelica, et
corporalis} (d.~2 c.~1), that the spiritual creature, the angelic, and the
corporeal were created together, and he sets aside the opinion of the ages
before the world with Hugh's words: \latin{salva tamen reverentia secretorum,
in quibus nihil temere asserendum est} (d.~2 c.~3). The angels were made in
the empyrean heaven, which was filled with angels as soon as it was made;
Lucifer's
``I will ascend into heaven'' means to the height of God (d.~2 cc.~4, 6;
Isa 14:13). They were made with four attributes: \latin{essentia simplex, id
est indivisibilis et immaterialis; et discretio personalis; et per rationem
naturaliter insitam, intelligentia, memoria et voluntas sive dilectio;
liberum quoque arbitrium} (d.~3 c.~1): a simple, indivisible, and
immaterial essence; personal distinctness; understanding, memory, and will or
love by the reason implanted in them by nature; and free choice. They
differed from the beginning in subtlety of essence, perspicacity of wisdom,
and readiness of choice, and those who excelled by nature received greater
gifts of grace (d.~3 c.~2), the text Aquinas quotes for the \emph{sed
contra} of \ST{I q.~62 a.~6}. Against those who held that the fallen were
created evil, the Master concludes: \latin{Ex praedictis igitur liquet,
Angelos omnes bonos esse creatos, et post creationem quosdam cecidisse a
bono quod habuissent, si perstitissent} (d.~3 c.~4): all the angels were
created good, and after their creation some fell from the good they would
have had if they had persevered, after a brief interval. They were neither
blessed nor miserable at their creation. Whether those who stood foreknew
their beatitude and so were in a manner blessed in hope, or were uncertain
of it like the rest, he answers: \latin{Mihi autem quod posterius dictum est
probabilius videtur} (d.~4): to me the latter seems more probable. They were
perfect in the first of Hugh's three perfections, not yet in the second
(d.~4).

\paragraph{Conversion and aversion; the fall (dd.~5--6).} In words he takes
from the Victorine \emph{Summa sententiarum}, the Master describes the
division:

\begin{quote}
\latin{Converti ad Deum fuit ei caritate adhaerere; averti, odio habere vel
invidere: invidiae namque mater est superbia, qua voluerunt se parificare
Deo. In conversis quasi in speculo relucere coepit Dei sapientia, qua
illuminati sunt; aversi vero excaecati sunt.} (d.~5 c.~1)
\end{quote}

\noindent To turn to God was to cleave to him in charity; to turn away was to
hate or envy him, and the mother of envy is pride, by which they willed to
make themselves equal to God. In those who turned, the wisdom of God began to
shine as in a mirror, and they were enlightened; those who turned away were
blinded. Those who stood were given cooperating grace to love God perfectly;
the fallen were not given it, and the fault was theirs, since they could
have stood until it was given (d.~5 cc.~3--5). Whether the good merited
their beatitude by the grace given them in confirmation, or received it as
a reward that their later service merits, the Master records both and
chooses the second: \latin{Et hoc magis mihi placere fateor} (d.~5 c.~6).
Some fell from the greater orders and some from the lesser, and among them
was one \latin{omnibus aliis cadentibus excellentior, nec inter stantes
aliquis eo fuit dignior} (d.~6 c.~1): more excellent than all who fell, and
none who stood was worthier. He is Job's ``beginning of the ways of God''
(Job 40:14) and Ezechiel's ``seal of resemblance'' covered with every
precious stone (Ezek 28:12--13), and he willed to be like the Most High
\latin{non per imitationem, sed per aequalitatem potentiae} (d.~6 c.~1), not
by imitation but by equality of power. The fallen were cast into this dark
air as into a prison until the judgment; some go down daily into hell with
the souls they torment; they rule over one another as the good do, some
over a province, some over a man, some over a vice (d.~6 cc.~2--5).

\paragraph{Confirmation, knowledge, and power (dd.~7--8).} The good cannot
sin and the wicked cannot will well, and yet both have free choice, the good
far freer than before (d.~7 cc.~1--3); the Master receives Isidore's
sentence, \latin{Angeli mutabiles natura, immutabiles sunt gratia} (d.~7
c.~4): the angels are changeable by nature and unchangeable by grace. The
demons keep a keen sense, knowing by subtlety of nature, long experience,
and the revelation of higher spirits (d.~7 c.~5); they are not creators,
nor are the good angels, though both can apply the hidden seeds of things
as God permits (d.~7 cc.~6--9). On the angels' bodies the Master sets out
the opinion that they have airy bodies and the opinion in which, he reports, many
catholic writers agree, that they are incorporeal and assume bodies for their ministry.
Of Augustine, who asks how the angels appeared to the Fathers and leaves it
open, he writes: \latin{Attende, lector, quia quaestionem propositam non
solvit, sed indiscussam relinquit} (d.~8 c.~3): mark, reader, that he does
not solve the question but leaves it undiscussed; and the Master holds
fast only that God never appeared to mortals in the form of his essence
(d.~8 c.~3). The demons do not enter the substance of the soul, which God
alone can fill, but work upon it (d.~8 c.~4, with Gennadius and Bede).

\paragraph{The orders (d.~9).} The ninth distinction gives the schools
their definition of an order and their account of the names. The Master
lists the orders as Hugh does, \latin{Angelos, Archangelos, Principatus,
Potestates, Virtutes, Dominationes, Thronos, Cherubim, Seraphim}, and then
arranges them in three triads, which he ascribes to Dionysius:
\latin{superiores: Seraphim, Cherubim, Throni; medii: Dominationes,
Principatus, Potestates; inferiores: Virtutes, Archangeli, Angeli} (d.~9
c.~1). The triads set the Principalities among the middle orders and the
Virtues among the lowest, which is Gregory's placing; the \emph{Celestial
Hierarchy} sets the Principalities in the lowest triad (\emph{CH} 9.1) and
the Virtues, Parker's ``Powers'', in the middle (\emph{CH} 8.1). An order is
defined as \latin{multitudo caelestium spirituum, qui inter se in aliquo
munere gratiae similantur, sicut et in naturalium datorum participatione
conveniunt} (d.~9 c.~2): a multitude of heavenly spirits who are alike in
some gift of grace, as they agree also in sharing natural gifts. Aquinas
takes this definition for the \emph{sed contra} of \ST{I q.~108 a.~4}. The
names are expounded in Gregory's words and descend in Gregory's sequence,
from the Seraphim, \latin{qui prae aliis ardent caritate}, who burn with
charity beyond the others, to the Angels, \latin{qui minora}, who announce
the lesser things (d.~9 c.~2). With Hugh the Master teaches
that \latin{Haec nomina illis non propter se, sed propter nos data sunt}
(d.~9 c.~3), and that each order is named from the gift it holds most
excellently, though all hold all; charity is a greater gift than knowledge,
and so the highest order is named from charity. The orders were not yet
distinguished before the fall, since the gifts in which they agree had not
yet been given; Scripture says that some fell from every order because some
of the fallen would have been in every order had they stood (d.~9 c.~4).
Within an order not all are equal (d.~9 c.~5). The ``tenth order'' made up
from men does not mean a tenth order of men, for Gregory teaches that men
are taken into the orders of the angels, some into the higher and some into
the lower, \latin{sed homines pro qualitate meritorum statuendos in ordinibus
Angelorum} (d.~9 c.~6): men are to be placed in the orders of the angels
according to the quality of their merits. Men are admitted to beatitude
according to the number not of those who fell but of those who stood, as
Gregory reads Deut 32:8 (d.~9 c.~7). Aquinas determines the same by grace
(\ST{I q.~108 a.~8}).

\paragraph{Missions and guardians (dd.~10--11).} Whether all the heavenly
spirits are sent, the Master reports the opinions Hugh gave, and records
that Michael, Gabriel, and Raphael are held by some to be of the highest
order, and that the names are either proper to single spirits or given to
whichever is sent for such a work (d.~10 cc.~1--2). The eleventh
distinction teaches the guardian angels:

\begin{quote}
\latin{Illud quoque sciendum est, quod Angeli boni deputati sunt ad
custodiam hominum, ita ut quisque electorum habeat Angelum, ad sui profectum
atque custodiam specialiter delegatum.} (d.~11 c.~1)
\end{quote}

\noindent The good angels are appointed to guard men, so that each of the
elect has an angel specially delegated for his progress and keeping. The
Master cites Jerome on Matthew 18:10, \latin{Magna dignitas animarum est, ut
unaquaeque habeat ab ortu nativitatis in sui custodiam delegatum Angelum}:
great is the dignity of souls, that each has from its birth an angel
delegated to guard it; and Gregory, that each man has a good angel for his
keeping and an evil angel for his testing (d.~11 c.~1). Aquinas takes
Jerome's sentence for the \emph{sed contra} of \ST{I q.~113 a.~2}. Whether
the angels grow in knowledge and beatitude until the judgment, the Master
records that Jerome and Augustine seem to disagree on the angels' knowledge
of the mystery hidden from the ages (Eph 3:9--10), reconciles them with
Haymo, and concludes: \latin{Constat itaque, omnes Angelos in cognitione
divinorum mysteriorum secundum processum temporis profecisse} (d.~11 c.~2):
all the angels have advanced with time in the knowledge of the divine
mysteries. Aquinas cites the distinction for the higher angels' long
knowledge of the Incarnation (\ST{I q.~106 a.~4 obj.~2}).

\sectionguard
\subsection{Bonaventure's \emph{Breviloquium}}\label{sec:fa-bonaventure}

Bonaventure, minister general of the Friars Minor and cardinal (d.~1274),
the Seraphic Doctor, summarised the faith in the \emph{Breviloquium}, and
its second part gives the angels three chapters: their production, the
apostasy of the demons, and the confirmation of the good (II.6--8).\footnote{The
Latin is that of the Quaracchi edition (1911).} His disputed positions on the
angels' matter and species are set out at \S\ref{sec:disputed}.

\paragraph{The production of the angels (II.6).} From the first moment the
angels have four attributes, which he takes from Hugh and the Master:
\latin{simplicitas essentiae, personalis discretio, propter rationem insitam
memoria, intelligentia et voluntas, et libertas arbitrii ad eligenda bona et
respuenda mala} (II.6.2): simplicity of essence, personal distinctness,
memory, understanding, and will by the reason implanted in them, and freedom
of choice to choose good and refuse evil. Four others accompany them:
\latin{virtuositas in operando, officiositas in ministrando, perspicacitas
in cognoscendo et immutabilitas post electionem sive in bono, sive in
malo} (II.6.2): power in acting, readiness in ministering, clearness in
knowing, and unchangeableness after their choice, whether in good or in
evil. The reason is the first Principle's own nature: having made all things
from nothing, it had to make not only what is near nothing, the bodily
creature, but also what is near itself, the intellectual substance, which is
most like God in simplicity and personal distinctness, bears the image of
the Trinity in memory, understanding, and will, and is free, so that it
might come by merit to glory (II.6.3).

\paragraph{The apostasy of the demons (II.7).} God made the angels good,
set between the highest Good and the changeable good of creatures:
turning to what is above, they would rise to grace and glory; turning to
what is below, they would fall into the evil of fault and punishment. Then:
\latin{Primus inter Angelos lucifer, praesumens de privato bono, privatam
appetiit excellentiam, volens aliis superferri; et ideo cecidit cum ceteris
consentientibus sibi} (II.7.1): the first among the angels, Lucifer,
presuming on his private good, sought a private excellence, wishing to be
raised above the others, and so fell with the rest who consented to him.
Bonaventure traces the fall to two acts: \latin{praesumendo constituit se
sibi principium, in se ipso gloriando; et ambiendo constituit se sibi
summum bonum, in se ipso quiescendo} (II.7.3): by presumption he made
himself his own principle, glorying in himself; by ambition he made himself
his own highest good, resting in himself. Since he was neither, he had to
fall by his disordered climb. Because the angels are unchangeable after
their choice, he was at once obstinate; and each of his four gifts was
turned: his will to hatred and envy of man, his clearness of reason to
deceptions, his readiness to minister to temptations, and his power,
diminished and restrained, to the working of wonders on bodies, so that he
seeks to be worshipped as God (II.7.4).

\paragraph{The confirmation of the good (II.8).} As those who turned away
were at once hardened, those who turned to God were at once confirmed by
grace and glory, perfectly enlightened, strengthened, and ordered, and this
in a threefold hierarchy:

\begin{quote}
\latin{Ad supremam autem spectant Throni, Cherubim et Seraphim; ad mediam
Dominationes, Virtutes et Potestates; ad infimam Principatus, Archangeli,
Angeli.} (II.8.1)
\end{quote}

\noindent To the highest hierarchy belong the Thrones, Cherubim, and
Seraphim; to the middle the Dominations, Virtues, and Powers; to the lowest
the Principalities, Archangels, and Angels. This is Dionysius's placing.
Many of them are sent in ministry and deputed to guard men, whom they
cleanse, enlighten, and perfect. Their contemplation is never disordered by
their ministry: since they see God face to face, \latin{quocumque mittantur,
intra Deum currunt} (II.8.2): wherever they are sent, they run within God,
as Gregory said before him. Bonaventure then derives the nine orders from
the four gifts made perfect in glory. Clearness of contemplation looks to
the divine majesty to be revered, to the truth to be understood, to the
goodness to be desired: hence the Thrones, the Cherubim, the Seraphim.
Perfect power is imperative, executive, and expeditive: hence the
Dominations, the Virtues, and the Powers, whose office is to drive back the
contrary powers. Perfect ministry is to rule, to reveal, and to raise up:
hence the Principalities, the Archangels, and the Angels, \latin{quia
custodiunt, ne stantes cadant, et cadentes adiuvant, ut resurgant} (II.8.3):
because they guard those who stand lest they fall, and help those who fall
to rise again. Every order takes its name from the gift it has received
most excellently (II.8.3), which is Gregory's rule; Aquinas explains the
names from the property of each order (\ST{I q.~108 a.~5};
\S\ref{sec:q108}).

\sectionguard
\subsection{The Medieval Doctors and Their Reception}

The teaching of the doctors of the eleventh to thirteenth centuries, with
the place where Aquinas or the Church receives it:

\begin{fourstable}{Witness}{Locus}{Teaching}{Received at}
Anselm & \emph{De casu diaboli} 3--4 & The devil did not receive perseverance because he would not; God did not give because he did not receive & \ST{I q.~62 a.~2}; \ST{I q.~63 a.~1}\\
Anselm & \emph{De casu diaboli} 4, 6 & The devil sinned by willing an advantage he did not have, against God's will, by a will subject to none; so he willed to be like God inordinately & \ST{I q.~63 a.~3} (Anselm named); \ST{I q.~63 a.~1}\\
Anselm & \emph{De casu diaboli} 12--14 & God gave the will for happiness and the will for justice together, so that justice might temper happiness without removing the power to exceed & \ST{I q.~59 a.~3}; \ST{I q.~60 a.~5}\\
Anselm & \emph{De casu diaboli} 6, 17, 24 & The good cannot desert what they held, as reward; the fallen cannot recover what they deserted, as penalty & \ST{I q.~62 a.~8}; \ST{I q.~64 a.~2}\\
Anselm & \emph{De casu diaboli} 27 & The evil will had no cause but itself & Lateran IV (DS~800); \ST{I q.~63 a.~1}\\
Anselm & \emph{Cur Deus homo} I.16--18 & The fallen are replaced from men; more men will be saved than angels fell; man was made for his own sake & \ST{I q.~63 a.~9}; \ST{I q.~108 a.~8}\\
Bernard & \emph{De consid.}\ V.4.7--8 & The angels' existence and blessedness are of faith, their intelligence of reason, the offices of their names of reverent conjecture; the nine offices in Gregory's order & \ST{I q.~108 aa.~5--6}\\
Bernard & \emph{De consid.}\ V.5.12 & The angel suggests good and does not create it; God alone is in the soul as its life & \ST{I q.~106 a.~2}; \ST{I q.~111 a.~2}\\
Bernard & \emph{In Cant.}\ 5; 19 & The angels need bodies only for ministry; each order loves God for what it contemplates & \ST{I q.~51 aa.~1--2}; \ST{I q.~108 a.~5}\\
Bernard & \emph{In Ps.\ Qui habitat} 11--12 & The angels ascend in contemplation and descend in compassion without losing the vision; reverence, devotion, and confidence toward the guardian angel & \ST{I q.~112 a.~1 ad~3}; \ST{I q.~113 a.~1} (s.c.\ Ps 90:11)\\
Bernard & \emph{In festo S.\ Michaelis} 1 & The angels love us for God's sake, for ours, and for their own, to rebuild their city & \ST{I q.~108 a.~8}\\
Hugh & \emph{De sacr.}\ I.5.4--8 & The angels were created together with corporeal matter, simple, distinct, wise, and free & \ST{I q.~61 a.~3}; Lateran IV (DS~800)\\
Hugh & \emph{De sacr.}\ I.5.24, 28 & Grace cooperates with those who turn and leaves those who turn away; God restrains the demons in four ways & \ST{I q.~62 a.~2}; \ST{I q.~109 a.~4}; \ST{I q.~114 a.~1}\\
Hugh & \emph{De sacr.}\ I.5.30--33 & Nine orders and man the tenth; more stood than fell; the names are given for our sake; three opinions on missions & \ST{I q.~63 a.~9}; \ST{I q.~108 a.~5}; \ST{I q.~112 a.~2}\\
Hugh & \emph{In Hier.\ cael.}\ VI & Love goes in where knowledge stays outside & \ST{I q.~108 a.~6 ad~3}\\
Honorius & \emph{Elucidarium} I.6--10 & A fixed number of the elect, nine orders and man the tenth; the first angel fell at once; the good were confirmed by their own merit & \ST{I q.~63 a.~6}; \ST{I q.~62 a.~8}; \ST{I q.~109 a.~4}\\
Honorius & \emph{Elucidarium} II.28 & Angels preside over nations and cities; each soul is entrusted to an angel; they appear in bodies of air & \ST{I q.~113 aa.~1--2}; \ST{I q.~51 a.~2}\\
Hildegard & \emph{Scivias} I.6 & The nine choirs in Gregory's sequence, each a lesson for the Church; they praise God for his works in the saints & \ST{I q.~108 a.~6}\\
Lombard & \emph{Sent.}\ II dd.~2--5 & Created together with the world, good, with four attributes; the stable were confirmed by cooperating grace & \ST{I q.~61 a.~3}; \ST{I q.~62 aa.~3, 6} (s.c.\ d.~3)\\
Lombard & \emph{Sent.}\ II d.~9 & Definition of an order; names from the more excellent gift; men placed in the orders by merit & \ST{I q.~108 a.~4} (s.c.); \ST{I q.~108 a.~8}\\
Lombard & \emph{Sent.}\ II d.~11 & Each of the elect has a guardian angel from birth; the angels advance in knowledge of the mysteries & \ST{I q.~113 a.~2} (s.c.\ Jerome); \ST{I q.~106 a.~4 obj.~2}\\
Bonaventure & \emph{Brev.}\ II.6--7 & Four attributes and four companions; Lucifer's presumption and ambition & \ST{I q.~63 aa.~2--3}\\
Bonaventure & \emph{Brev.}\ II.8 & The Dionysian triads; the orders derived from the four gifts; wherever sent, they run within God & \ST{I q.~108 a.~6}; \ST{I q.~112 a.~1 ad~3}\\
\end{fourstable}
